% ========================================================
%  习近平新时代中国特色社会主义思想概论  期末全真模拟试卷（二）（含参考答案）
%  覆盖范围：导论、第一至八章、第十至十二章 核心考点
% ========================================================
\documentclass[12pt, a4paper]{article}

% ---- 中文支持 ----
\usepackage[UTF8]{ctex}

% ---- 数学 ----
\usepackage{amsmath, amssymb, amsthm}
\usepackage{mathtools}

% ---- 页面布局 ----
\usepackage[top=2.5cm, bottom=2.5cm, left=2.8cm, right=2.8cm]{geometry}

% ---- 表格 ----
\usepackage{booktabs, array, multirow, makecell}
\usepackage{tabularx}

% ---- 页眉页脚 ----
\usepackage{fancyhdr}
\usepackage{lastpage}
\setlength{\headheight}{14pt}
\pagestyle{fancy}
\fancyhf{}
\renewcommand{\headrulewidth}{0.6pt}
\lhead{\small 习近平新时代中国特色社会主义思想概论·Claude·B卷}
\rhead{\small 共 \pageref{LastPage} 页}
\cfoot{\small 第 \thepage\ 页}

% ---- 计数器与样式 ----
\usepackage{enumitem}
\usepackage{xcolor}
\usepackage{tcolorbox}
\tcbuselibrary{skins, breakable}

% ---- 填空线 ----
\newcommand{\blank}[1]{\underline{\hspace{#1}}}

% ---- 选择题选项列表 ----
\newlist{choices}{enumerate}{1}
\setlist[choices]{label=\textbf{\Alph*.}, itemsep=1pt, topsep=2pt, parsep=0pt, leftmargin=2.4em}

% ---- 答案盒子 ----
\newtcolorbox{answerbox}[1][]{
  breakable,
  colback=gray!6,
  colframe=gray!50,
  fonttitle=\bfseries\small,
  title=参考答案,
  #1
}

% ---- 题号格式 ----
\newcounter{bigq}
\newcommand{\BigQ}[1]{%
  \stepcounter{bigq}%
  \vspace{6pt}%
  \noindent{\large\bfseries 第\chinese{bigq}大题\quad #1}%
  \vspace{4pt}\par
}

\newcounter{subq}[bigq]
\newcommand{\SubQ}{%
  \stepcounter{subq}%
  \noindent\textbf{（\arabic{subq}）}\quad
}

% 分隔线
\newcommand{\examrule}{\noindent\rule{\linewidth}{0.6pt}}

% ---- 文档开始 ----
\begin{document}

% ============================================================
%  封面
% ============================================================
\begin{center}
  {\LARGE\bfseries 习近平新时代中国特色社会主义思想概论}\\[8pt]
  {\large\bfseries 期末全真模拟测试卷（二）}\\[16pt]
  \begin{tabular}{lll}
    考试时间：90 分钟 & \quad & 满分：100 分 \\[4pt]
    \multicolumn{3}{l}{注意事项：请将答案写在答题纸指定区域，写在本卷上无效。}
  \end{tabular}
  \vspace{4pt}

  \examrule

  \vspace{6pt}
  \begin{tabular}{|c|c|c|c|c|}
    \hline
    \textbf{题号} & 一（单选） & 二（多选） & 三（简答） & \textbf{总分} \\
    \hline
    \textbf{满分} & 30 & 20 & 50 & 100 \\
    \hline
    \textbf{得分} & & & & \\
    \hline
  \end{tabular}
\end{center}

\vspace{10pt}
\examrule
\vspace{6pt}

% ============================================================
%  第一大题：单项选择题
% ============================================================
\BigQ{单项选择题（每题只有一个正确答案，每题 1 分，共 30 分）}

\SubQ 关于"两个结合"，下列表述正确的是？
\begin{choices}
\item 马克思主义基本原理同中国具体实际相结合，同中华优秀传统文化相结合
\item 中国具体实际同中华优秀传统文化相结合
\item 科学社会主义同资本主义相结合
\item 马克思主义同自由主义相结合
\end{choices}

\SubQ "十四个坚持"在思想体系中的作用是？
\begin{choices}
\item 这一思想的世界观和方法论
\item 新时代坚持和发展中国特色社会主义的基本方略
\item 全面展示了这一思想的实践成果
\item 这一思想最为核心关键的组成部分
\end{choices}

\SubQ "十三个方面成就"指的是？
\begin{choices}
\item 这一思想的基本方略
\item 全面展示了这一思想的实践成果
\item 这一思想的世界观和方法论
\item 核心关键的组成部分
\end{choices}

\SubQ "六个必须坚持"体现的是？
\begin{choices}
\item 实践成果
\item 基本方略
\item 世界观和方法论
\item 核心关键组成部分
\end{choices}

\SubQ 关于"两个确立"的决定性意义，下列表述不准确的是？
\begin{choices}
\item 是新时代取得一切成就的根本原因
\item 是战胜一切艰难险阻的最大确定性、最大底气、最大保证
\item 确立了习近平同志党中央的核心、全党的核心地位
\item 是马克思主义中国化时代化"第三次历史性飞跃"的标志
\end{choices}

\SubQ 科学社会主义基本原则在哪部著作及其手稿中充分展开并深入论证？
\begin{choices}
\item 《共产党宣言》
\item 《资本论》
\item 《哥达纲领批判》
\item 《反杜林论》
\end{choices}

\SubQ 道路自信是对什么的自信？
\begin{choices}
\item 中国特色社会主义制度具有制度优势的自信
\item 发展方向和未来命运的自信
\item 中国特色社会主义文化先进性的自信
\item 马克思主义理论科学性、真理性的自信
\end{choices}

\SubQ 党的基本路线被称为？
\begin{choices}
\item 党的生命线、人民的幸福线
\item 国家的生命线、人民的幸福线
\item 党和国家的根本制度
\item 社会主义初级阶段的总纲领
\end{choices}

\SubQ "五位一体"总体布局不包括？
\begin{choices}
\item 经济建设
\item 政治建设
\item 国防建设
\item 生态文明建设
\end{choices}

\SubQ 中国梦的本质是？
\begin{choices}
\item 国家富强、民族振兴、人民幸福
\item 经济发展、社会稳定、文化繁荣
\item 科技进步、教育发达、人才济济
\item 改革开放、创新驱动、协调发展
\end{choices}

\SubQ 中国式现代化的本质要求不包括？
\begin{choices}
\item 坚持中国共产党领导
\item 发展全过程人民民主
\item 优先发展军事工业
\item 促进人与自然和谐共生
\end{choices}

\SubQ 中国式现代化创造人类文明新形态的表现不包括？
\begin{choices}
\item 深深植根于中华优秀传统文化
\item 体现科学社会主义的先进本质
\item 完全排斥借鉴人类优秀文明成果
\item 代表人类文明进步的发展方向
\end{choices}

\SubQ 推进中国式现代化需要正确处理的重大关系中，"效率与公平的关系"主要强调？
\begin{choices}
\item 效率优先、不顾公平
\item 在发展中兼顾效率与公平，二者不可偏废
\item 公平优先、放弃效率
\item 效率与公平互不相关
\end{choices}

\SubQ 党的领导制度是？
\begin{choices}
\item 我国的根本领导制度
\item 一般性的工作制度
\item 临时性的协调机制
\item 地方性管理制度
\end{choices}

\SubQ 维护党中央权威和集中统一领导，是什么的重大建党原则？
\begin{choices}
\item 一般政党
\item 成熟的马克思主义执政党
\item 地方党组织
\item 群众组织
\end{choices}

\SubQ "事在四方，要在中央"，强调的是？
\begin{choices}
\item 地方自主权的扩大
\item 党中央的核心领导地位和定于一尊的权威
\item 基层民主的重要性
\item 党员个人的自由裁量权
\end{choices}

\SubQ "民心是最大的政治"，这一论断强调的是？
\begin{choices}
\item 民心决定事业兴衰成败
\item 民心与政治无关
\item 民心只是次要因素
\item 民心可以被忽视
\end{choices}

\SubQ 调查研究是？
\begin{choices}
\item 脱离群众路线的独立方法
\item 获得真知灼见的源头活水，是贯彻群众路线的有效途径
\item 一种过时的工作方式
\item 只适用于经济领域
\end{choices}

\SubQ 全面深化改革要坚持的正确方法论不包括？
\begin{choices}
\item 加强顶层设计和摸着石头过河相结合
\item 统筹改革发展稳定
\item 坚持重大改革于法有据
\item 改革可以脱离法治单独推进
\end{choices}

\SubQ "六个紧紧围绕"路线图中，深化经济体制改革紧紧围绕的是？
\begin{choices}
\item 建设社会主义核心价值体系
\item 使市场在资源配置中起决定性作用和更好发挥政府作用
\item 提高科学执政、民主执政、依法执政水平
\item 更好保障和改善民生
\end{choices}

\SubQ 高质量发展，就是从什么转向什么？
\begin{choices}
\item 从"有没有"转向"好不好"
\item 从"快不快"转向"稳不稳"
\item 从"大不大"转向"强不强"
\item 从"新不新"转向"旧不旧"
\end{choices}

\SubQ 构建新发展格局，最本质的特征是？
\begin{choices}
\item 扩大对外贸易规模
\item 实现高水平的自立自强
\item 提高进口依存度
\item 扩大产能过剩行业规模
\end{choices}

\SubQ 关键核心技术受制于人是？
\begin{choices}
\item 我们最大的隐患
\item 可以忽略的小问题
\item 只影响个别企业的局部问题
\item 与国家安全无关的问题
\end{choices}

\SubQ 人民代表大会制度是？
\begin{choices}
\item 我国的根本政治制度
\item 我国的基本政治制度之一
\item 临时性政治安排
\item 地方性政治制度
\end{choices}

\SubQ 全过程人民民主是全链条、全方位、全覆盖的民主，其中"全链条"主要体现在哪些环节？
\begin{choices}
\item 民主选举、民主协商、民主决策、民主管理、民主监督
\item 仅指民主选举环节
\item 仅指民主监督环节
\item 仅指民主决策环节
\end{choices}

\SubQ 新闻舆论工作做好的第一位要求是？
\begin{choices}
\item 提高传播速度
\item 坚持正确的政治方向
\item 扩大覆盖范围
\item 增加娱乐内容
\end{choices}

\SubQ 社会主义核心价值观在国家层面的内容是？
\begin{choices}
\item 自由、平等、公正、法治
\item 爱国、敬业、诚信、友善
\item 富强、民主、文明、和谐
\item 团结、互助、友爱、进步
\end{choices}

\SubQ "枫桥经验""浦江经验"主要体现在？
\begin{choices}
\item 城乡社区治理、社会矛盾纠纷多元预防调处化解
\item 国防建设
\item 对外贸易
\item 文化产业发展
\end{choices}

\SubQ 生态文明制度体系建设中提到的"三条红线"是？
\begin{choices}
\item 生态保护红线、环境质量底线、资源利用上线
\item 财政红线、金融红线、贸易红线
\item 教育红线、科技红线、人才红线
\item 民生红线、安全红线、稳定红线
\end{choices}

\SubQ 中国参与全球环境治理坚持的原则不包括？
\begin{choices}
\item 坚持以人为本
\item 坚持科学治理
\item 坚持单边主义
\item 坚持共同但有区别的责任原则
\end{choices}

\vspace{10pt}
\examrule

% ============================================================
%  第二大题：多项选择题
% ============================================================
\BigQ{多项选择题（每题有两个或两个以上正确答案，多选、少选、错选均不得分；每题 2 分，共 20 分）}

\SubQ 关于"两个结合"，下列说法正确的是？
\begin{choices}
\item 坚持把马克思主义基本原理同中国具体实际相结合
\item 坚持把马克思主义基本原理同中华优秀传统文化相结合
\item 马克思主义基本原理是"魂脉"
\item 中华优秀传统文化是"根脉"
\item 中国具体实际是"魂脉"
\end{choices}

\SubQ 习近平新时代中国特色社会主义思想的历史地位包括？
\begin{choices}
\item 当代中国马克思主义、二十一世纪马克思主义
\item 中华文化和中国精神的时代精华
\item 实现了马克思主义中国化时代化新的飞跃
\item 是党和国家必须长期坚持的指导思想
\item 实现了马克思主义中国化时代化"第三次历史性飞跃"
\end{choices}

\SubQ "五位一体"总体布局包括？
\begin{choices}
\item 经济建设
\item 政治建设
\item 文化建设
\item 社会建设
\item 生态文明建设
\end{choices}

\SubQ "四个全面"战略布局包括？
\begin{choices}
\item 全面建设社会主义现代化国家
\item 全面深化改革
\item 全面依法治国
\item 全面从严治党
\item 全面推进军事现代化
\end{choices}

\SubQ 中国式现代化的本质要求包括？
\begin{choices}
\item 坚持中国共产党领导
\item 坚持中国特色社会主义
\item 实现高质量发展
\item 发展全过程人民民主
\item 优先发展重工业
\end{choices}

\SubQ 中国特色社会主义制度和国家治理体系的显著优势包括？
\begin{choices}
\item 坚持党的集中统一领导
\item 坚持人民当家作主
\item 坚持全面依法治国
\item 坚持全国一盘棋，集中力量办大事
\item 坚持地方分权自治、各自为政
\end{choices}

\SubQ 全面深化改革开放的正确方法论包括？
\begin{choices}
\item 增强全面深化改革的系统性、整体性、协同性
\item 加强顶层设计和摸着石头过河相结合
\item 统筹改革发展稳定
\item 坚持重大改革于法有据
\item 改革可以脱离法治单独推进，提高效率
\end{choices}

\SubQ 社会主义基本经济制度的新概括包括？
\begin{choices}
\item 公有制为主体、多种所有制经济共同发展
\item 按劳分配为主体、多种分配方式并存
\item 社会主义市场经济体制
\item 完全计划经济体制
\item 完全私有化经济体制
\end{choices}

\SubQ 全过程人民民主是最广泛、最真实、最管用的民主，体现在？
\begin{choices}
\item 全体人民共同持续参与，各民族共同平等享有
\item 真实反映人民的期盼、希望和诉求
\item 保证在党的领导下有效治理国家
\item 巩固平等团结互助和谐的社会主义民族关系
\item 仅在选举投票环节起作用
\end{choices}

\SubQ 我国社会主义现代化建设的基础性、战略性支撑包括？
\begin{choices}
\item 教育
\item 科技
\item 人才
\item 军事
\item 外交
\end{choices}

\vspace{10pt}
\examrule

% ============================================================
%  第三大题：简答题
% ============================================================
\BigQ{简答题（请简要论述，只写出结论不进行论述不能得满分；每题 10 分，共 50 分）}

\SubQ（10 分）"两个确立"的决定性意义是什么？

\vspace{6pt}
\SubQ（10 分）如何理解新时代的内涵？为什么说中国特色社会主义进入新时代是我国发展新的历史方位？

\vspace{6pt}
\SubQ（10 分）如何理解人民民主是社会主义的生命？

\vspace{6pt}
\SubQ（10 分）为什么说文化繁荣兴盛是实现中华民族伟大复兴的必然要求？

\vspace{6pt}
\SubQ（10 分）为什么必须把生态文明建设摆在全局工作的突出位置？

\vspace{16pt}
\examrule
\vspace{4pt}
\begin{center}
  {\large\bfseries ——试题到此结束，请认真检查——}
\end{center}
\vspace{4pt}
\examrule

% ============================================================
%  参考答案
% ============================================================
\newpage
\setcounter{bigq}{0}

\begin{center}
  {\LARGE\bfseries 参考答案与评分标准}
\end{center}
\vspace{8pt}
\examrule

% -------- 第一大题参考答案 --------
\BigQ{单项选择题参考答案}

\begin{answerbox}

\textbf{答案速查：}

\begin{center}
\begin{tabular}{|c|c|c|c|c|c|c|c|c|}
\hline
题号 & 1 & 2 & 3 & 4 & 5 & 6 & 7 & 8 \\
\hline
答案 & A & B & B & C & D & B & B & B \\
\hline
\end{tabular}

\vspace{4pt}

\begin{tabular}{|c|c|c|c|c|c|c|c|c|}
\hline
题号 & 9 & 10 & 11 & 12 & 13 & 14 & 15 & 16 \\
\hline
答案 & C & A & C & C & B & A & B & B \\
\hline
\end{tabular}

\vspace{4pt}

\begin{tabular}{|c|c|c|c|c|c|c|c|c|}
\hline
题号 & 17 & 18 & 19 & 20 & 21 & 22 & 23 & 24 \\
\hline
答案 & A & B & D & B & A & B & A & A \\
\hline
\end{tabular}

\vspace{4pt}

\begin{tabular}{|c|c|c|c|c|c|c|}
\hline
题号 & 25 & 26 & 27 & 28 & 29 & 30 \\
\hline
答案 & A & B & C & A & A & C \\
\hline
\end{tabular}
\end{center}

\vspace{8pt}
\textbf{简要解析：}

\begin{enumerate}[itemsep=2pt, topsep=4pt, leftmargin=1.6em]
\item "两个结合"指马克思主义基本原理同中国具体实际相结合、同中华优秀传统文化相结合，B、C、D均为错误搭配。
\item "十四个坚持"（十四个方略）是新时代坚持和发展中国特色社会主义的基本方略，注意与"十个明确""六个必须坚持"区分。
\item "十三个方面成就"全面展示了这一思想的实践成果，是科学体系四个组成部分之一。
\item "六个必须坚持"是这一思想的世界观和方法论，彰显了理论品格与鲜明特征。
\item D项为干扰项：课堂明确强调思想的历史地位只讲"新的飞跃"，没有官方"第三次历史性飞跃"提法，"两个确立"本身也不包含这一表述。
\item 科学社会主义基本原则首次在《共产党宣言》中全面阐述，在《资本论》及其手稿中充分展开并深入论证，在《哥达纲领批判》《反杜林论》中进一步阐发，三部著作分工不同需区分记忆。
\item 道路自信是对发展方向和未来命运的自信；理论自信对应科学性真理性，制度自信对应制度优势，文化自信对应文化先进性，四者不可混淆。
\item 党的基本路线是国家的生命线、人民的幸福线，是社会主义初级阶段的基本路线，概括为"一个中心、两个基本点"。
\item "五位一体"总体布局是经济建设、政治建设、文化建设、社会建设、生态文明建设，不包括国防建设，国防和军队建设是另一独立战略领域。
\item 中国梦的本质是国家富强、民族振兴、人民幸福，三个维度辩证统一。
\item 中国式现代化的本质要求共九个方面，均未涉及"优先发展军事工业"，C项为典型干扰项。
\item 中国式现代化创造人类文明新形态恰恰强调"借鉴吸收一切人类优秀文明成果"，C项"完全排斥"与原文表述截然相反。
\item 效率与公平的关系，强调的是在发展中统筹兼顾、不可偏废，而非片面强调任何一方。
\item 党的领导制度是我国的根本领导制度，是中国特色社会主义制度体系的核心。
\item 维护党中央权威和集中统一领导，是一个成熟的马克思主义执政党的重大建党原则，体现党的成熟度和组织原则。
\item "事在四方，要在中央"强调党中央是大脑和中枢，必须有定于一尊、一锤定音的权威。
\item "民心是最大的政治"强调民心决定事业兴衰成败，是第四章的重要论断。
\item 调查研究是获得真知灼见的源头活水，是贯彻群众路线的有效途径，"没有调查就没有发言权，就没有决策权"。
\item 改革和法治如鸟之两翼、车之两轮，必须坚持重大改革于法有据，D项"脱离法治单独推进"违背这一方法论。
\item "六个紧紧围绕"中，深化经济体制改革紧紧围绕使市场在资源配置中起决定性作用和更好发挥政府作用，注意区分政治、文化、社会、生态、党建体制改革各自围绕的内容。
\item 高质量发展就是从"有没有"转向"好不好"，体现发展理念由量的扩张转向质的提升。
\item 构建新发展格局最本质的特征是实现高水平的自立自强，这是着力加快科技自立自强的核心目标。
\item 核心技术受制于人是我们最大的隐患，不掌握核心技术就会被"卡脖子""牵鼻子"。
\item 人民代表大会制度是我国的根本政治制度，注意与基本政治制度（多党合作和政治协商制度、民族区域自治制度、基层群众自治制度）层级区分。
\item 全过程人民民主"全链条"体现在民主选举、民主协商、民主决策、民主管理、民主监督五个环节，并非只局限于某一环节。
\item 新闻舆论工作有鲜明的意识形态属性，坚持正确的政治方向是做好新闻舆论工作的第一位要求。
\item 社会主义核心价值观国家层面内容是富强、民主、文明、和谐；社会层面是自由、平等、公正、法治；个人层面是爱国、敬业、诚信、友善，三个层面不能混淆。
\item "枫桥经验""浦江经验"主要体现在城乡社区治理和社会矛盾纠纷多元预防调处化解机制方面，是加强党组织领导的基层治理体系的重要内容。
\item 生态文明制度体系建设要加快划定并严守生态保护红线、环境质量底线、资源利用上线三条红线。
\item 中国全球环境治理坚持以人为本、科学治理、多边主义、共同但有区别的责任原则，C项"单边主义"与中国一贯坚持的多边主义立场相悖，属典型干扰项。
\end{enumerate}

\end{answerbox}

% -------- 第二大题参考答案 --------
\BigQ{多项选择题参考答案}

\begin{answerbox}

\textbf{答案速查：}

\begin{center}
\begin{tabular}{|c|c|c|c|c|c|}
\hline
题号 & 1 & 2 & 3 & 4 & 5 \\
\hline
答案 & ABCD & ABCD & ABCDE & ABCD & ABCD \\
\hline
\end{tabular}

\vspace{4pt}

\begin{tabular}{|c|c|c|c|c|c|}
\hline
题号 & 6 & 7 & 8 & 9 & 10 \\
\hline
答案 & ABCD & ABCD & ABC & ABCD & ABC \\
\hline
\end{tabular}
\end{center}

\vspace{8pt}
\textbf{简要解析：}

\begin{enumerate}[itemsep=2pt, topsep=4pt, leftmargin=1.6em]
\item E项为干扰项：马克思主义基本原理是"魂脉"，中华优秀传统文化是"根脉"，"中国具体实际"不是"两个结合"中"魂脉"或"根脉"的指代对象，而是相结合的另一方。
\item E项为干扰项：思想的历史地位只讲"新的飞跃"，没有官方"第三次历史性飞跃"的表述。
\item "五位一体"总体布局五项内容均正确，应全选。
\item E项"全面推进军事现代化"为干扰项，"四个全面"战略布局准确表述为全面建设社会主义现代化国家、全面深化改革、全面依法治国、全面从严治党。
\item E项为干扰项，中国式现代化本质要求原文共九个方面，未提及"优先发展重工业"这一表述。
\item E项"地方分权自治、各自为政"为干扰项，与"坚持全国一盘棋、集中力量办大事"的显著优势直接相悖。
\item E项为干扰项，改革和法治相辅相成、相伴而生，不能脱离法治单独推进。
\item D、E项为干扰项，社会主义基本经济制度新概括是公有制为主体多种所有制经济共同发展、按劳分配为主体多种分配方式并存、社会主义市场经济体制三者，不是完全计划经济或完全私有化。
\item E项为干扰项，全过程人民民主是全链条、全方位、全覆盖的民主，并不仅限于选举投票环节。
\item D、E项为干扰项，原文明确强调教育、科技、人才三者是全面建设社会主义现代化国家的基础性、战略性支撑，未将军事、外交列入其中。
\end{enumerate}

\end{answerbox}

% -------- 第三大题参考答案 --------
\BigQ{简答题参考答案}

\begin{answerbox}

\textbf{第（1）题　"两个确立"的决定性意义是什么？（10 分）}

"两个确立"的内容是：确立习近平同志党中央的核心、全党的核心地位，确立习近平新时代中国特色社会主义思想的指导地位。

"两个确立"的决定性意义主要体现在：第一，"两个确立"是新时代党和人民事业取得历史性成就、发生历史性变革的根本原因；第二，"两个确立"是战胜一切艰难险阻、应对一切不确定性的最大确定性、最大底气、最大保证，为新征程上应对各种风险挑战提供了最根本的政治保证。

需要注意辨析的是，"两个确立"与"两个维护"紧密联系但内涵不同："两个确立"侧重于政治原则和理论指导地位的确立，"两个维护"的关键则是维护党中央权威和集中统一领导，落到实处是维护党中央和各级党组织的领导地位，二者应结合起来理解，不能混为一谈。

必须把坚定拥护"两个确立"的决定性意义，转化为坚定捍卫习近平总书记党中央的核心、全党的核心地位，自觉做到"两个维护"的实际行动。

\textit{评分标准："两个确立"内容阐述2分；决定性意义展开4分；与"两个维护"的联系区别辨析2分；总结2分。}

\vspace{6pt}
\textbf{第（2）题　如何理解新时代的内涵？为什么说中国特色社会主义进入新时代是我国发展新的历史方位？（10 分）}

中国特色社会主义新时代的内涵包括五个方面：第一，是承前启后、继往开来，在新的历史条件下继续夺取中国特色社会主义伟大胜利的时代；第二，是决胜全面建成小康社会、进而全面建设社会主义现代化强国的时代；第三，是全国各族人民团结奋斗、不断创造美好生活、逐步实现全体人民共同富裕的时代；第四，是全体中华儿女勠力同心、奋力实现中华民族伟大复兴中国梦的时代；第五，是我国不断为人类作出更大贡献的时代。

之所以说这是我国发展新的历史方位，是因为以上五层含义从历史延续性、发展目标、人民生活、民族复兴、世界贡献五个维度，全面刻画了中国特色社会主义所处的历史阶段和发展坐标，标志着我国发展站在了新的历史起点上，是关系全局的历史性变化。

需要强调的是，新时代的判断是以社会主要矛盾变化为基础作出的，但这一判断并未改变我国仍处于并将长期处于社会主义初级阶段的基本国情，也没有改变我国是世界最大发展中国家的国际地位，即"一个变、两个没有变"。

\textit{评分标准：五层内涵每点1.2分（共6分）；"新的历史方位"的理解2分；结合"一个变两个没有变"加以辨析2分。}

\vspace{6pt}
\textbf{第（3）题　如何理解人民民主是社会主义的生命？（10 分）}

理解人民民主是社会主义的生命，应把握以下几个层面：

第一，民主具有普遍性与特殊性的统一。民主是全人类共同价值，是人类政治文明发展的成果；但民主同时是历史的、具体的，一个国家实行什么样的民主制度，必须与这个国家的社会历史条件和现实国情相适应，实现民主的道路并非只有一条。

第二，人民民主具有鲜明的社会主义性质。人民民主建立在社会主义经济基础之上，体现了社会主义国家的性质，反映了社会主义制度的本质要求，是一种新型的社会主义民主。

第三，人民民主具有重要地位。人民民主是全面建设社会主义现代化国家的应有之义，发展社会主义民主政治是其内在要求和重要目标，社会主义现代化是全体人民的现代化，必然要求具有高度发达的民主政治。

第四，实践要求上，必须坚持党的领导、人民当家作主、依法治国有机统一，积极发展全过程人民民主，不断扩大人民有序政治参与，保障人民当家作主。

总之，没有民主就没有社会主义，更没有社会主义现代化；人民民主是社会主义的生命，是检验和彰显社会主义制度优越性的重要标尺。

\textit{评分标准：民主的普遍性与特殊性2分；人民民主的社会主义性质2分；人民民主的重要地位2分；实践要求2分；总结2分。}

\vspace{6pt}
\textbf{第（4）题　为什么说文化繁荣兴盛是实现中华民族伟大复兴的必然要求？（10 分）}

文化繁荣兴盛是实现中华民族伟大复兴的必然要求，主要理由有四个方面：第一，文化繁荣兴盛是实现中华民族伟大复兴的精神支撑，没有文化的繁荣兴盛就没有民族复兴的强大精神力量；第二，文化繁荣兴盛是建设社会主义现代化强国的应有之义，社会主义现代化国家应当是物质文明和精神文明协调发展的国家；第三，文化繁荣兴盛是满足人民日益增长的美好生活需要的内在要求，人民的美好生活需要呈现多样化、多层次、多方面特点，精神文化需求是其重要组成部分；第四，文化繁荣兴盛是在世界文化激荡中站稳脚跟的前提基础，唯有文化自立自强才能在世界文化竞争中保持自身特色和立场。

坚定文化自信是文化繁荣兴盛的基础：文化自信是更基础、更广泛、更深沉的力量，其中中华优秀传统文化是坚定文化自信的深厚基础，党在伟大斗争中孕育的革命文化和社会主义先进文化是坚强基石，中国特色社会主义伟大实践是现实基础。

为此必须坚持中国特色社会主义文化发展道路，坚持举旗帜、聚民心、育新人、兴文化、展形象，坚持"二为"方向、"双百"方针，最终铸就社会主义文化新辉煌，为民族复兴提供坚实精神支撑。

\textit{评分标准：四个方面理由每点1.5分（共6分）；文化自信的基础地位2分；实践要求与总结2分。}

\vspace{6pt}
\textbf{第（5）题　为什么必须把生态文明建设摆在全局工作的突出位置？（10 分）}

必须把生态文明建设摆在全局工作的突出位置，主要原因如下：

第一，从文明兴衰规律看，生态兴则文明兴，生态衰则文明衰，生态环境是人类生存和发展的根基，生态环境变化直接影响文明兴衰演替。

第二，从现实紧迫性看，改革开放以来我国经济发展取得巨大成就，但也积累了大量生态环境问题，必须清醒认识保护生态环境、治理环境污染的紧迫性和艰巨性，坚定不移走生产发展、生活富裕、生态良好的文明发展道路。

第三，从民生意义看，良好的生态环境是最公平的公共产品，是最普惠的民生福祉；保护生态环境就是保护自然价值和增值自然资本，就是保护经济社会发展潜力和后劲，能够使绿水青山持续发挥生态效益和经济社会效益。

第四，从制度安排看，党的十八大把生态文明建设纳入中国特色社会主义"五位一体"总体布局，党中央提出一系列新理念新思想新战略，形成了习近平生态文明思想，并要求加快划定并严守生态保护红线、环境质量底线、资源利用上线三条红线。

总之，必须坚持人与自然和谐共生，把生态文明建设贯穿于经济社会发展全过程和各方面，为人民创造良好生产生活环境。

\textit{评分标准：文明兴衰规律2分；现实紧迫性2分；民生意义2分；制度安排2分；总结2分。}

\end{answerbox}

\vspace{12pt}
\examrule
\begin{center}
  \textit{参考答案仅供参考，评分时应酌情给分，论述题言之有理、要点齐全且能展开论述可得相应分数。}
\end{center}

\end{document}
